\subsection{Predictive benchmarks}
\label{sec:nash-selector}

A key limitation of our statistical framework is that comparing the predictive power of different AI simulations provides no absolute benchmark for how well these samples predict human responses in general. 
Thus, we require a suitable theoretical or statistical benchmark for a more comprehensive analysis.
However, due to the scale and heterogeneity of our games, several appealing benchmarks are impractical.

Ideally, we would apply the standard level-$k$ model from Section~\ref{sec:predict_new_AR} to generate predictive distributions across these games. 
Unfortunately, no existing mechanical method reliably identifies which choices correspond to specific levels of reasoning across such diverse contexts. 
For example, the ``obvious'' choice for a level-$0$ player is not always the highest number---particularly when bonuses are large or the bonus rule involves selecting number 11. 
Consequently, higher-level reasoning does not follow the intuitive progression featured by the original game in \citeauthor{1120Arad2012}.

Another seemingly attractive alternative would be to follow the approaches of \cite{FudenbergLiang2019PredictPlay}, \cite{hirasawa2022using}, or \cite{andrews2025transfer}, which involve training a bespoke supervised machine learning model on past game data to predict responses. 
However, this method is infeasible for our purposes because of two problems. 
First, it relies heavily on having access to a large and representative dataset, which we currently do not possess.\footnote{The combined approach of \cite{zhu2025capturing}---using machine learning to flexibly estimate parameters in a well-defined behavioral model and then using that model to predict responses---is infeasible for the same reason.} 
Second, as discussed in Section~\ref{sec:approach}, traditional machine learning models (e.g., regression, decision trees, or generic neural networks) are not well-suited for predicting human behavior in novel settings because they cannot adapt to structural differences across settings.
Even with substantial representative data, if a game had even one new parameter value (e.g., a new bonus rule or an upper bound above the highest from the training data), the model would need to be retrained from scratch.

We employ two complementary benchmarks that each address a different part of the above limitations.
First, we use the cognitive hierarchy model of \cite{Camerer2004Cognitive}.
Throughout the economics literature, it is one of the most empirically accurate predictors of human responses in one-shot strategic games.
It can also produce predictions for all the structurally distinct games in $S$, but does require some prior information, which we discuss in the next subsection.
Second, we use Nash equilibrium with Harsanyi-Selten selection, which requires no empirical calibration, but, along with most equilibrium concepts, has not been particularly successful at predicting initial human play for static games in the literature. 

The challenge of finding an ideal benchmark actually highlights a unique strength of \aisub{s}: they can generate plausibly accurate predictions in virtually any setting without extensive training data or pre-specified behavioral parameters.
To our knowledge, no other well-known benchmarks are both likely to have predictive accuracy and flexible enough to be applied to such a wide range of structurally distinct games.
